\documentclass[12pt]{article}
\newcommand{\DocShort}{Ajuste exponencial con funciones de Excel}
\input{preamble_population.tex}
\begin{document}
\doccover{Ajuste exponencial por mínimos cuadrados con funciones de Excel}{2026-05}
\handoutintro{Ajustar un modelo exponencial con mínimos cuadrados en Excel y expresarlo de forma físicamente interpretable.}{Series con crecimiento multiplicativo sostenido.}

\section{Datos ejemplo}
Aquí se usa $t$ como tiempo medido en años desde 1970.

\begin{center}
\begin{tabular}{lcccccc}
\toprule
Año & 1970 & 1980 & 1990 & 2000 & 2010 & 2020 \\
\midrule
$t$ (años) & 0 & 10 & 20 & 30 & 40 & 50 \\
Población & 18,500 & 23,100 & 28,700 & 35,600 & 44,300 & 54,800 \\
\bottomrule
\end{tabular}
\end{center}

\section{Ecuaciones mínimas}
La forma que devuelve Excel puede escribirse como:
\begin{equation}
P=A\,B^t
\end{equation}
donde $P$ es la población proyectada, $A$ es una constante, $B$ es el factor de crecimiento por año y $t$ es el tiempo en años desde el origen elegido.

Como
\begin{equation}
B=e^r
\end{equation}
también es posible escribir el modelo como:
\begin{equation}
P=Ae^{rt}
\end{equation}
donde $r$ es la tasa continua de crecimiento anual.

En Excel:
\[
\texttt{=LOGEST(rango\_P, rango\_t, TRUE, TRUE)}
\]
En Excel en español:
\[
\texttt{=ESTIMACION.LOGARITMICA(rango\_P; rango\_t; VERDADERO; VERDADERO)}
\]

\excelnote{En \texttt{LOGEST}, la variable dependiente es la población. Su equivalente en Excel en español es \texttt{ESTIMACION.LOGARITMICA}. Luego conviene calcular $r=\ln(B)$ para expresar el modelo en la forma $P=Ae^{rt}$.

\textbf{Nota.} Las funciones \texttt{GROWTH} y \texttt{CRECIMIENTO} son útiles para proyectar valores, pero no devuelven directamente los parámetros de la ecuación.}

\section{Procedimiento}
\begin{enumerate}
\item Construir la variable temporal $t$.
\item Ingresar $t$ y $P$ en Excel.
\item Aplicar \texttt{LOGEST}.
\item Calcular $r=\ln(B)$.
\item Escribir el modelo preferentemente como $P=Ae^{rt}$.
\item Evaluar el modelo para los años futuros.
\end{enumerate}

\section{Ejemplo resuelto}
Con los datos dados, \texttt{LOGEST} entrega aproximadamente:
\[
A=\num{18556.81},\qquad B=1.021947
\]

Entonces,
\[
r=\ln(1.021947)=0.02171\ \text{año}^{-1}
\]

La forma preferida del modelo es:
\[
P=\num{18556.81}\,e^{0.02171\,t}
\]

Para los años futuros:
\[
P_{2030}=\num{18556.81}\,e^{0.02171(60)}=\num{68268}
\]
\[
P_{2040}=\num{18556.81}\,e^{0.02171(70)}=\num{84820}
\]
\[
P_{2050}=\num{18556.81}\,e^{0.02171(80)}=\num{105387}
\]

\resultbox{El ajuste exponencial con LOGEST produce 68,268 habitantes en 2030, 84,820 en 2040 y 105,387 en 2050.}

\section{Teoría breve}
El modelo exponencial supone que la tasa relativa de crecimiento es aproximadamente constante. Expresar el resultado mediante la tasa continua $r$ es útil porque esa constante tiene significado físico directo como tasa anual de crecimiento.

\section{Observaciones de uso}
\begin{itemize}
\item Suele dar proyecciones más altas que un ajuste lineal.
\item Puede sobreestimar horizontes largos si la ciudad se acerca a saturación.
\item Debe aclararse siempre la unidad temporal de $t$ y, por tanto, la unidad inversa de $r$.
\end{itemize}
\end{document}
